%=====================================================================
\chapter{What Virtualization Is}\label{ch:what}
%=====================================================================
\locator{1}{Part I --- Foundations of Virtualization}

\begin{outcomes}
By the end of this chapter you will be able to:
\begin{tight}
  \item \textbf{Define} virtualization as a layer that presents a resource
        different from the one it sits on, and \textbf{identify} that layer in
        a system you are shown.
  \item \textbf{Distinguish} a virtual machine from a process, an emulator and
        a simulator, by what each one is allowed to assume about the hardware.
  \item \textbf{Explain} the four reasons an organisation inserts such a layer
        --- consolidation, isolation, encapsulation and elasticity --- and
        \textbf{state} which of the four motivated each milestone in the
        field's history.
  \item \textbf{Compute} a consolidation ratio and the power it saves, from
        measured utilisation figures.
  \item \textbf{Describe} what is being virtualized in a given product, and
        locate the chapter of this book that covers it.
\end{tight}
\end{outcomes}

\begin{prereq}
Nothing in this book. You need the idea that an operating system stands between
a program and the hardware, and that a program cannot tell the difference
between a real device and a convincing imitation of one. If the second half of
that sentence sounds improbable, this chapter is the right place to start.
\end{prereq}

\begin{keyterms}
virtualization $\bullet$ virtual machine $\bullet$ hypervisor $\bullet$ virtual
machine monitor $\bullet$ host $\bullet$ guest $\bullet$ abstraction $\bullet$
multiplexing $\bullet$ aggregation $\bullet$ emulation $\bullet$ consolidation
$\bullet$ consolidation ratio $\bullet$ utilisation $\bullet$ server sprawl
$\bullet$ isolation $\bullet$ encapsulation $\bullet$ elasticity $\bullet$
time-sharing
\end{keyterms}

\section{A Room Full of Idle Servers}

Walk into a university machine room that has not been consolidated and you will
find something strange: forty servers, each costing as much as a car, each
running at eight per cent of its capacity. The library catalogue has a server.
The examination system has a server. The staff portal has a server. None of them
is busy. Together they draw fourteen kilowatts and need as much again in cooling,
and the department pays for all of it.

Nobody planned this. It arrived one purchase order at a time, and each purchase
was defensible. The examination system must not be slowed down by the library
catalogue, so it gets its own machine. The portal needs a newer version of the
database than the catalogue can tolerate, so it gets its own machine. The
vendor's support contract is void if anything else runs on the box, so it gets
its own machine. Every reason is a reason about \emph{isolation}, and the only
isolation mechanism available was a second computer.

This is \term{server sprawl}, and it is the problem virtualization was brought
to industry to solve. The insight is that the forty applications genuinely do
need to be isolated from one another --- that part was never wrong --- but that
isolation does not have to be purchased one chassis at a time.

\begin{definitionbox}[Virtualization]
\term{Virtualization} is the insertion of a software layer that presents a
resource which differs from the one physically present, while preserving the
interface that software above the layer already expects.

\smallskip
The layer does one of three things, or several at once:
\begin{tight}
  \item \term{multiplexing} --- one physical resource is presented as many
        (one processor as twelve vCPUs);
  \item \term{aggregation} --- many physical resources are presented as one
        (twelve disks as one volume);
  \item \term{abstraction} --- a resource is presented as a different kind of
        thing entirely (a file on a disk presented as a disk).
\end{tight}

\smallskip
What makes it virtualization rather than merely software is the second clause:
the interface is \emph{preserved}. An unmodified operating system, written with
no knowledge that any of this is happening, runs correctly above the layer.
\end{definitionbox}

\dsfig{%
\begin{tikzpicture}[font=\small,
  hw/.style={rounded corners=2pt, draw=neutral, line width=1pt, fill=neutral!12,
             text=neutral!40!black, minimum width=11.6cm, minimum height=0.95cm,
             align=center, font=\small\bfseries},
  vl/.style={rounded corners=2pt, draw=primary, line width=1.4pt, fill=primary!14,
             text=primary, minimum width=11.6cm, minimum height=0.95cm,
             align=center, font=\small\bfseries},
  g/.style={rounded corners=2pt, draw=secondary, line width=1pt, fill=secondary!8,
            text=secondary!80!black, minimum width=3.5cm, minimum height=1.45cm,
            align=center, font=\scriptsize},
  % align=left is load-bearing: a \\ in a node whose style sets neither align
  % nor text width is "Not allowed in LR mode".
  lab/.style={font=\scriptsize\itshape, text=neutral, anchor=west, align=left}
]
  \node[hw] (hw) at (0,0) {Physical hardware \quad
        {\scriptsize\mdseries 64 cores \ $\bullet$\ 512 GB \ $\bullet$\ 2 NICs \ $\bullet$\ 8 disks}};
  \node[vl] (vl) at (0,1.25) {Virtualization layer};
  \node[g]  (g1) at (-4.05,2.95) {\textbf{Guest A}\\[2pt]
        4 vCPU \ $\bullet$\ 16 GB\\1 NIC \ $\bullet$\ 1 disk};
  \node[g]  (g2) at (0,2.95)     {\textbf{Guest B}\\[2pt]
        2 vCPU \ $\bullet$\ 8 GB\\1 NIC \ $\bullet$\ 1 disk};
  \node[g]  (g3) at (4.05,2.95)  {\textbf{Guest C}\\[2pt]
        8 vCPU \ $\bullet$\ 64 GB\\2 NIC \ $\bullet$\ 3 disks};

  \foreach \n in {g1,g2,g3}{
    \draw[-{Stealth[length=2mm]}, primary!55, line width=0.9pt] (\n) -- (\n |- vl.north);}
  \draw[-{Stealth[length=2mm]}, primary!55, line width=0.9pt] (vl) -- (hw);

  \node[lab] at (6.1,2.95) {each believes it owns\\a whole computer};
  \node[lab] at (6.1,1.25) {maintains the belief};
  \node[lab] at (6.1,0)    {what actually exists};

  \node[font=\scriptsize, text=accent, anchor=east] at (-6.1,2.95)
       {unmodified};
  \node[font=\scriptsize, text=accent, anchor=east] at (-6.1,2.55)
       {operating systems};
\end{tikzpicture}}%
{The layer, and the only claim it makes. Each guest is told a different machine
exists, and none of the three was modified to be told. Every later chapter
replaces one row of this diagram with something more specific.}{the-layer}

\begin{definitionbox}[Virtual Machine, Hypervisor, Host, Guest]
A \term{virtual machine} is the illusion the layer sells: a complete computer
--- processor, memory, firmware, disks, network interfaces --- implemented in
software, on which an operating system can be installed as if on metal.

\smallskip
The \term{hypervisor}, equivalently the \term{virtual machine monitor} or VMM,
is the software that creates and maintains that illusion. The \term{host} is the
physical machine and the software running directly on it; a \term{guest} is an
operating system running inside a virtual machine.

\smallskip
This book uses \emph{hypervisor} and \emph{VMM} interchangeably, as the
literature does. \chref{vmm} is the one place the distinction is sharpened,
because that is where it begins to matter.
\end{definitionbox}

\section{What It Is Not}

Three neighbouring ideas are routinely confused with virtualization, and the
confusion is worth clearing before it compounds.

\begin{conceptbox}[Virtualization, Emulation, Simulation]
\term{Emulation} interprets every guest instruction in software, translating
from one instruction set to another. An ARM emulator on an x86 laptop reads an
ARM instruction, works out what it means, and performs the equivalent x86 work.
It is slow --- commonly ten to a hundred times slower than native --- and
completely general: the guest architecture need not resemble the host at all.

\smallskip
\term{Virtualization} runs guest instructions \emph{on the real processor}
wherever it safely can, and intervenes only for the instructions it cannot allow
to execute directly. It is therefore fast --- commonly within a few per cent of
native --- and restricted: the guest must be built for the same architecture as
the host.

\smallskip
\emph{Simulation} models a system's behaviour without claiming to run its
software: a network simulator predicts throughput; it does not boot a router.

\smallskip
QEMU does the first two in different modes, which is why it appears on both
sides of the comparison in \chref{architectures}. Running
\texttt{qemu-system-aarch64} on an x86 host is emulation. Running
\texttt{qemu-system-x86\_64} with KVM enabled is virtualization, and the same
binary does both.
\end{conceptbox}

\begin{alertbox}[A container is not a virtual machine, and both are virtualization]
The two most common statements about containers in a first course contradict
each other, and both are half right.

\smallskip
A container does not contain a virtual machine: there is no virtual processor,
no virtual firmware, no guest kernel. Processes in a container run on the host's
kernel as ordinary processes, and \texttt{ps} on the host can see them.

\smallskip
But a container \emph{is} virtualization by the definition above. The layer is
inserted inside the kernel rather than beneath it, and what it multiplexes is
the operating system's own namespace of process identifiers, mount points,
network interfaces and users. The interface preserved is the system-call
interface, not the instruction set. \chref{containers} makes this precise; the
reason it matters this early is that ``virtual machine or container?'' is a
question about \emph{which layer}, not about whether a layer exists.
\end{alertbox}

\section{Why Anyone Would Do This}

Four motives, and they arrived in this order historically. Every product in this
book was built for at least one of them, and the ones that lasted serve several.

\rowcolors{2}{white}{lightbg}
\begin{longtable}{@{}L{2.9cm}L{6.3cm}L{5.8cm}@{}}
\toprule
\rowcolor{primary!15}
\textbf{Motive} & \textbf{What it buys} & \textbf{Where this book shows it} \\
\midrule
\endhead
\term{Consolidation} & Many lightly-loaded workloads on one well-loaded
machine. Fewer machines, less power, less cooling, less floor space, less to
patch & The worked example below; \chref{cpumem} on how far it can be pushed;
\chref{performance} on when it has gone too far \\
\term{Isolation} & A fault, a crash or a compromise in one workload does not
reach the others. This is the reason the forty servers existed, and
virtualization preserves it while removing the forty & \chref{vmm} on what makes
isolation sound; \chref{security} on where it fails \\
\term{Encapsulation} & A running machine becomes a file. It can be copied,
versioned, moved, snapshotted and restored --- operations no physical server
supports & \chref{storage} on the file; \chref{migration} on moving it while it
runs \\
\term{Elasticity} & Capacity becomes a quantity you request rather than a
chassis you procure. Minutes instead of months, and returnable & \chref{cloud}
entire; \chref{kubernetes} for the same idea one layer up \\
\bottomrule
\end{longtable}

\begin{conceptbox}[Encapsulation is the one with no physical analogue]
Consolidation is an efficiency argument and isolation is a safety argument, and
a sufficiently rich organisation could buy its way out of both. Encapsulation is
different in kind.

\smallskip
Once a running computer is a file, every operation you know how to perform on a
file becomes available to it. You can take a snapshot before a risky upgrade and
return to it in four seconds when the upgrade fails. You can copy a production
machine into a test environment, bit for bit, including the fault you are trying
to reproduce. You can keep last Tuesday. None of this has a physical equivalent
--- there is no way to take a snapshot of a server in a rack --- and it is the
capability administrators miss most when they go back.
\end{conceptbox}

\begin{worked}[The Forty Servers, Costed]
The machine room of the opening section, measured over one semester. Decide how
many hosts replace it, and what that saves.

\smallskip
\textbf{What is there.} 40 servers, each 2 sockets $\times$ 8 cores $=$ 16
cores and 64 GB of memory. Measured per server: mean CPU utilisation 8\%, peak
22\%; mean memory in use 11 GB. Each draws 350 W.

\smallskip
\textbf{Step 1 --- what is actually consumed.}
\begin{tight}
  \item CPU: $40 \times 16 \times 0.08 = 51.2$ cores' worth of work, spread over
        640 physical cores.
  \item Memory: $40 \times 11 = 440$ GB in use, of 2{,}560 GB installed.
\end{tight}

\smallskip
\textbf{Step 2 --- size the replacement on memory, not CPU.} Memory is the
binding constraint in almost every consolidation, because it cannot be
oversubscribed as freely as processor time (\chref{cpumem} says why). Allow 30\%
headroom for growth and for the hypervisor itself:
\[ 440 \times 1.3 = 572 \text{ GB} \]
On hosts of 512 GB that is $\lceil 572/512 \rceil = 2$ hosts. Add one more so
that a single host failure does not take the estate down --- an \emph{N+1}
design --- giving \textbf{3 hosts}.

\smallskip
\textbf{Step 3 --- check that CPU is not the real constraint.} Give each
workload 4 vCPUs, which is generous against a measured peak of
$16 \times 0.22 = 3.5$ cores. Three hosts of 2 $\times$ 32 cores give 192
physical cores:
\[ \sigma = \frac{40 \times 4}{192} = \frac{160}{192} = 0.83 \]
An overcommit ratio below 1 means that even if every guest demanded its full
allocation simultaneously, the cores exist. CPU is not the constraint; memory
was right.

\smallskip
\textbf{Step 4 --- the consolidation ratio.}
\[ R = \frac{40 \text{ workloads}}{3 \text{ hosts}} = \mathbf{13.3 : 1} \]

\smallskip
\textbf{Step 5 --- the power.} The new hosts are larger and draw 700 W each.
\[ 40 \times 350 = 14{,}000 \text{ W} \qquad 3 \times 700 = 2{,}100 \text{ W}
   \qquad \text{saved: } 11.9 \text{ kW} \]
Cooling roughly tracks the load; at a power usage effectiveness of 1.6 the
facility saves $11.9 \times 1.6 = 19.0$ kW. Over a year:
\[ 19.0 \times 8{,}760 = 166{,}400 \text{ kWh} \]
At PKR 45 per kWh that is \textbf{about PKR 7.5 million a year}, before counting
the thirty-seven servers not bought, the rack space not rented, or the
thirty-seven operating systems not patched.

\smallskip
\textbf{What the number does not say.} Those three hosts are now a shared fate.
The arithmetic above bought a 13:1 reduction in machines and a 13:1
concentration of risk, and \chref{migration} is where that bill comes due.
\end{worked}

\dsfig{%
\begin{tikzpicture}[font=\scriptsize]
  % --- before: 40 boxes, lightly used ---
  \node[font=\small\bfseries, text=neutral!40!black, anchor=south west]
       at (0,2.7) {Before: 40 servers at 8\%};
  \foreach \i in {0,...,39}{
    \pgfmathsetmacro{\cx}{mod(\i,10)*0.62}
    \pgfmathsetmacro{\cy}{2.35-floor(\i/10)*0.56}
    \draw[draw=neutral!55, line width=0.6pt, fill=white]
          (\cx,\cy) rectangle ++(0.50,0.44);
    \fill[neutral!45] (\cx,\cy) rectangle ++(0.50,0.0352);}
  \node[font=\scriptsize\itshape, text=neutral, anchor=north west] at (0,0.28)
       {640 cores installed, 51 used \ $\bullet$\ 14 kW};

  % --- after: 3 boxes, well used ---
  \node[font=\small\bfseries, text=primary, anchor=south west]
       at (7.6,2.7) {After: 3 hosts at 68\%};
  \foreach \i in {0,1,2}{
    \pgfmathsetmacro{\cx}{7.6+\i*1.45}
    \pgfmathsetmacro{\mx}{7.6+\i*1.45+0.575}
    \draw[draw=primary, line width=1pt, fill=white]
          (\cx,1.23) rectangle ++(1.15,1.56);
    \fill[primary!35] (\cx,1.23) rectangle ++(1.15,1.0608);
    \node[font=\scriptsize, text=primary!70!black] at (\mx,0.98)
         {13--14 guests};}
  \node[font=\scriptsize\itshape, text=primary, anchor=north west] at (7.6,0.28)
       {192 cores installed, 160 allocated \ $\bullet$\ 2.1 kW};

  \draw[-{Stealth[length=2.6mm]}, primary!55, line width=1.1pt]
        (6.55,1.55) -- (7.35,1.55);
  \node[font=\scriptsize\bfseries, text=accent] at (6.95,1.85) {13.3 : 1};

  \node[font=\scriptsize\itshape, text=neutral!50!black, anchor=north west]
       at (0,-0.15)
       {Shaded area is the fraction of installed capacity actually doing work.};
\end{tikzpicture}}%
{The same work, before and after. Nothing was made faster; the idle capacity was
simply stopped from being bought. The right-hand picture is also the risk
picture: three failure domains where there were forty.}{consolidation}

\section{Where It Came From}

Virtualization is not a 2000s invention. It is a 1960s invention that the
personal computer made irrelevant for thirty years and the data centre made
urgent again.

\dsfig{%
% The colour of each event is written out per node rather than carried by a
% parameterised ev=<colour> style: a TikZ style taking #1 does not survive
% being an argument of \dsfig.  See the README's figure rules.
\begin{tikzpicture}[font=\scriptsize,
  ev/.style={rounded corners=2pt, line width=0.8pt, font=\scriptsize,
             align=center, inner sep=3pt, anchor=south}
]
  \draw[-{Stealth[length=2.6mm]}, neutral!60, line width=1.2pt]
        (-0.4,0) -- (16.2,0);
  \foreach \x/\y in {0/1967, 2.1/1974, 4.0/1999, 5.6/2003, 7.1/2005,
                     8.7/2006, 10.6/2013, 12.0/2014, 13.6/2018, 15.3/2021}{
    \draw[neutral!60, line width=0.9pt] (\x,-0.12) -- (\x,0.12);
    \node[font=\scriptsize\bfseries, text=neutral!40!black, anchor=north]
         at (\x,-0.16) {\y};}

  \node[ev, draw=primary, fill=primary!10, text=primary!60!black]
       at (0,0.28)    {IBM CP-67\\\emph{first hypervisor}};
  \node[ev, draw=primary, fill=primary!10, text=primary!60!black]
       at (2.1,1.42)  {Popek \&\\Goldberg\\\emph{the theory}};
  % Every \\ is at the top level of the node.  A \\ inside the \emph{} group
  % is "Undefined control sequence" at the end of the figure, because an
  % aligned node redefines \\ only at its own brace level.
  \node[ev, draw=secondary, fill=secondary!10, text=secondary!60!black]
       at (4.0,0.28)  {VMware\\\emph{x86, by}\\\emph{translation}};
  \node[ev, draw=secondary, fill=secondary!10, text=secondary!60!black]
       at (5.6,1.42)  {Xen\\\emph{x86, by}\\\emph{asking nicely}};
  \node[ev, draw=secondary, fill=secondary!10, text=secondary!60!black]
       at (7.1,0.28)  {Intel VT-x\\\emph{x86, in}\\\emph{silicon}};
  \node[ev, draw=purplehl, fill=purplehl!10, text=purplehl!60!black]
       at (8.7,1.42)  {Amazon EC2\\\emph{sold by}\\\emph{the hour}};
  \node[ev, draw=greenhl, fill=greenhl!10, text=greenhl!60!black]
       at (10.6,0.28) {Docker\\\emph{the image,}\\\emph{not the VM}};
  \node[ev, draw=greenhl, fill=greenhl!10, text=greenhl!60!black]
       at (12.0,1.42) {Kubernetes\\\emph{a thousand}\\\emph{of them}};
  \node[ev, draw=greenhl, fill=greenhl!10, text=greenhl!60!black]
       at (13.6,0.28) {Firecracker\\\emph{the VM}\\\emph{comes back}};
  \node[ev, draw=goldhl, fill=goldhl!10, text=goldhl!60!black]
       at (15.3,1.42) {SEV-SNP, TDX\\\emph{the guest stops}\\\emph{trusting the host}};

  \draw[decorate, decoration={brace, amplitude=4pt, mirror}, accent, line width=0.8pt]
        (2.4,-0.62) -- (3.7,-0.62);
  \node[font=\scriptsize\itshape, text=accent, anchor=north] at (3.05,-0.76)
       {25 years in which x86 could not do this};
\end{tikzpicture}}%
{The field in ten dates, coloured by the part of this book that covers each.
The gap in the middle is the subject of \chref{vmm}: the theory was settled in
1974, and the dominant architecture then spent a quarter of a century unable to
satisfy it.}{timeline}

IBM built CP-67 for the System/360 Model 67 because its customers had a problem
of scarcity rather than sprawl. A mainframe cost millions and there was exactly
one; a dozen departments each needed their own operating environment on it, and
several of those operating systems could not be persuaded to share. CP-67 gave
each department a complete virtual System/360 of its own. Its successor VM/370
ran in production for decades, and its descendant z/VM still does.

Then the minicomputer and the personal computer made computers cheap, and the
problem CP-67 solved stopped being a problem: if you needed a second computing
environment you bought a second computer. The idea went quiet for twenty-five
years, kept alive on mainframes and in research, while the architecture that won
the market --- x86 --- turned out, for reasons \chref{vmm} sets out exactly, to
be incapable of supporting it correctly.

\begin{conceptbox}[Why the second wave looked nothing like the first]
CP-67 was built for scarcity: one machine, many users. VMware was built for the
opposite --- forty machines, each almost idle. The motive inverted completely,
and the mechanism was reused unchanged.

\smallskip
That is worth noticing because the same inversion happened again. Containers
were adopted in 2013 largely for \emph{packaging}: a reproducible image that
runs the same on a laptop and in production. Only afterwards did the industry
discover that the same mechanism gave it density --- a thousand containers on a
host that would hold fifty virtual machines --- and density became the reason
people cite today. A mechanism outlives the motive that produced it, which is
why this book teaches mechanisms.
\end{conceptbox}

\section{What Is Virtualized, and Where to Find It}

``Virtualization'' names a technique, not a product category. The useful
question about any system is \emph{which resource} it virtualizes, because that
determines what it can promise.

\rowcolors{2}{white}{lightbg}
\begin{longtable}{@{}L{3.3cm}L{5.0cm}L{3.0cm}L{3.7cm}@{}}
\toprule
\rowcolor{primary!15}
\textbf{What is virtualized} & \textbf{So that} & \textbf{Examples} &
\textbf{Chapter} \\
\midrule
\endhead
The whole machine & An unmodified operating system boots on it &
ESXi, Hyper-V, KVM, Xen & \chref{platforms} \\
The operating system & Processes see a private system, sharing one kernel &
Docker, containerd, LXC & \chref{containers}, \chref{docker} \\
The processor and memory & Several guests divide cores and RAM & vCPU
scheduling, ballooning & \chref{cpumem} \\
A device & A guest drives hardware it does not exclusively own & virtio,
SR-IOV, VFIO & \chref{paravirt} \\
Storage & A file behaves as a disk; many disks behave as one & qcow2, VMDK,
Ceph, vSAN & \chref{storage} \\
The network & A switch exists in software; a LAN spans data centres & Open
vSwitch, VXLAN, SDN & \chref{network} \\
The desktop & The screen is here and the computer is elsewhere & VDI, Citrix,
AVD & \chref{desktop} \\
An accelerator & One GPU serves several tenants & vGPU, MIG, SR-IOV &
\chref{gpu} \\
\bottomrule
\end{longtable}

\begin{pitfall}
\begin{tight}
  \item \textbf{Treating ``virtual'' as meaning ``not real''.} A virtual
        processor executes real instructions on real silicon and consumes real
        electricity. What is virtual is the \emph{mapping}, not the work. The
        fastest way to mislead yourself about capacity is to believe the vCPU
        count is free.
  \item \textbf{Sizing a consolidation from installed capacity rather than
        measured use.} The forty servers above had 2{,}560 GB installed and used
        440 GB. Planning against the first number buys five times the hardware
        you need; it is the single most common and most expensive mistake in
        this subject.
  \item \textbf{Assuming consolidation is free.} It converts capital cost into
        shared risk and contention. Three hosts instead of forty is also three
        failure domains instead of forty, and every guest on a host is now a
        potential noisy neighbour to every other (\chref{performance}).
  \item \textbf{Calling a hypervisor an operating system, or vice versa.} They
        schedule different things for different consumers. The confusion makes
        \chref{vmm} incomprehensible, because the whole chapter turns on what
        the hypervisor is allowed to assume that an operating system is not.
  \item \textbf{Believing the marketing distinction between ``virtualization''
        and ``cloud''.} A cloud is virtualization with a billing system, an API
        and somebody else's staff. \chref{cloud} is explicit about which of the
        three is doing the work in any given claim.
\end{tight}
\end{pitfall}

\begin{lab}[Your First Virtual Machine]
Runs on any Linux host, or on Windows under WSL2 with VirtualBox substituted at
step 2. Appendix~A installs the prerequisites. Expect fifteen minutes.

\begin{lstlisting}[language=bash]
#!/usr/bin/env bash
# ch01_first_vm.sh -- create, snapshot, break and restore a virtual machine.
set -euo pipefail

# --- 1. Does this processor support virtualization at all? --------------
# vmx = Intel VT-x, svm = AMD-V. Chapter 4 explains what these flags buy.
# A count of 0 means the feature is absent or disabled in firmware.
grep -c -E '\<(vmx|svm)\>' /proc/cpuinfo
lscpu | grep -i 'virtuali[sz]ation' || echo "no virtualization line reported"

# --- 2. Install the hypervisor and confirm it is usable -----------------
sudo apt-get install -y qemu-kvm libvirt-daemon-system virtinst \
                        libguestfs-tools cloud-image-utils
# virt-host-validate prints PASS/WARN/FAIL for every prerequisite.
sudo virt-host-validate qemu | head -20

# --- 3. Fetch a cloud image and give it a login -------------------------
# A cloud image is a disk file with an operating system already installed:
# encapsulation (section 1.3) before you have even started.
wget -nc -O ubuntu.img \
  https://cloud-images.ubuntu.com/jammy/current/jammy-server-cloudimg-amd64.img
cat > user-data <<'YAML'
#cloud-config
password: vss2026
chpasswd: { expire: False }
ssh_pwauth: True
YAML
cloud-localds seed.iso user-data

# --- 4. Create the virtual machine --------------------------------------
# Read the arguments as a parts list for an imaginary computer.
sudo cp ubuntu.img /var/lib/libvirt/images/vss1.qcow2
sudo qemu-img resize /var/lib/libvirt/images/vss1.qcow2 10G
sudo virt-install --name vss1 --memory 2048 --vcpus 2 \
  --disk /var/lib/libvirt/images/vss1.qcow2,format=qcow2 \
  --disk seed.iso,device=cdrom --os-variant ubuntu22.04 \
  --import --graphics none --noautoconsole

# --- 5. Look at it from the host ----------------------------------------
# dominfo is the guest as the hypervisor sees it: an allocation, not a machine.
sudo virsh list --all
sudo virsh dominfo vss1

# --- 6. Snapshot it, then break it on purpose ---------------------------
sudo virsh snapshot-create-as vss1 clean "before I broke it"
sudo virsh console vss1    # log in as ubuntu / vss2026, then:
                           #   sudo rm -rf /etc/systemd  ; sudo reboot
                           # Ctrl-] returns you to the host.

# --- 7. Undo four minutes of damage in four seconds ---------------------
# No physical server supports this operation. That is section 1.3's point.
sudo virsh snapshot-revert vss1 clean --running
sudo virsh snapshot-list vss1

# --- 8. Tear down -------------------------------------------------------
sudo virsh destroy vss1 || true
sudo virsh undefine vss1 --remove-all-storage --snapshots-metadata
\end{lstlisting}

\textbf{What to notice.} Step 6 took longer to type than step 7 took to undo,
and step 7 has no physical equivalent. Everything else in this book is detail;
that asymmetry is the subject.
\end{lab}

\begin{checkpoint}
\begin{tightnum}
  \item A department runs 24 servers, each with 128 GB of memory, of which a
        measured 19 GB is in use. You are sizing hosts with 384 GB each, with
        30\% headroom and an N+1 spare. How many hosts, and what consolidation
        ratio does that give?
  \item A colleague says ``containers are faster than virtual machines because
        there is no virtualization layer.'' Which half of that sentence is
        wrong, and what is the correct statement?
  \item Of the four motives in Section 1.3, which one has no equivalent for
        physical servers at all? Give an operation it makes possible and say
        why a physical machine cannot offer it.
\end{tightnum}
\end{checkpoint}

\begin{chaptersummary}
\begin{tight}
  \item Virtualization inserts a layer that presents a resource differing from
        the physical one, while preserving the interface above it. It
        multiplexes, aggregates or abstracts --- often all three.
  \item The preserved interface is what distinguishes it: an unmodified
        operating system runs above the layer.
  \item Emulation interprets every instruction and is general and slow;
        virtualization executes most instructions natively and is restricted
        and fast. QEMU does both, in different modes.
  \item Four motives: consolidation, isolation, encapsulation, elasticity.
        Encapsulation --- a running machine as a file --- is the one with no
        physical analogue.
  \item Size a consolidation from measured use, not installed capacity, and
        expect memory rather than CPU to be the binding constraint.
  \item The technique is from 1967. The twenty-five-year gap in the middle of
        its history is an x86 problem, and \chref{vmm} is what went wrong.
  \item Ask of any product: \emph{which resource does it virtualize?} That
        determines what it can promise, and which chapter of this book covers
        it.
\end{tight}
\end{chaptersummary}

\begin{reviewq}
  \item \textbf{(MCQ)} The defining property that separates virtualization from
        ordinary resource-sharing software is that: (a)~it runs in kernel mode;
        (b)~it preserves the interface, so unmodified software runs above it;
        (c)~it improves performance; (d)~it requires hardware support.
  \item \textbf{(MCQ)} A tool runs ARM binaries on an x86 laptop at roughly one
        fiftieth of native speed. This is: (a)~virtualization;
        (b)~paravirtualization; (c)~emulation; (d)~containerisation.
  \item \textbf{(Short)} Give the three things a virtualization layer can do to
        a resource, with one concrete example of each from a system you have
        used.
  \item \textbf{(Short)} Explain why memory, rather than processor time, is
        usually the constraint that decides how many guests fit on a host.
  \item \textbf{(Short)} The opening section argued that forty separate servers
        was a rational response to a real requirement. Name the requirement,
        and explain precisely what virtualization changed about it.
  \item \textbf{(Applied)} A faculty runs 60 servers at a measured mean of 6\%
        CPU and 14 GB of 96 GB memory each, drawing 300 W apiece. Size a
        replacement estate on 768 GB hosts drawing 850 W, with 30\% headroom and
        N+1. Give the host count, the consolidation ratio and the annual energy
        saved at PUE 1.6 and PKR 45 per kWh. Then state the one risk your
        proposal has introduced and what you would do about it.
  \item \textbf{(Applied)} For each of these, say which resource is virtualized
        and which chapter of this book covers it: a laptop running Windows
        inside a window on macOS; a Kubernetes cluster; an NVMe disk presented
        to eight tenants; an office where staff keep no local files and log in
        to the same desktop from any machine; a GPU shared between three
        training jobs.
\end{reviewq}
